\section{Experiments}\label{sec:experiments}
\subsection{Datasets and Few-Shot Splits}
We evaluate MVTec AD \cite{mvtec2019} and VisA \cite{visa2022}. The strict external-transfer audit additionally incorporates MPDD. MVTec AD comprises industrial object and texture categories with pixel-level defect masks. VisA contains a broader range of industrial categories and serves as the primary benchmark for full-scale corruption and calibration shift. MPDD provides six metal-part categories for an external stress test that transfers from MVTec to MPDD; within-MPDD category certification is not pursued because fewer than six non-target categories remain for a three-way source split. For each category, support sets are sampled from the normal training partition using fixed seeds. The CRESS source archive is separate: it consists of held-out known-normal evaluation images from non-target source categories (the \texttt{test/good} partition for MVTec and the corresponding normal partition for VisA and MPDD). Source normal status is therefore used, but target test labels never enter reference construction, threshold proposal, or certification. Unless otherwise specified, $k\in\{1,2,4,8\}$ is employed for the clean accuracy-storage and strict nested studies, whereas the target-only corruption tables use $k\in\{4,8\}$. The $k=1$ strict cells use the declared patch-split support calibration fallback and are interpreted as stress tests; LOIO calibration itself requires $k\geq2$.

\subsection{Evaluation Protocols}
The target-only corruption audit evaluates all test images from all 15 MVTec and 12 VisA categories at $k\in\{4,8\}$ under five nested support seeds and four corruptions. Its primary p-values follow Eq.~\eqref{eq:loio}: calibration scores use LOIO fits and each test score uses the full-support fit. The fold-matched construction is reported only as a sensitivity analysis. Alarm comparisons use the predeclared tolerance $10^{-6}$ to absorb float32 representation error at attainable grid points.

The primary empirical groups are the target-only corruption audit and the frozen strict nested CRESS audit; the category-count feasibility curves are analytic design calculations. Supporting groups cover clean ranking and storage, PCA64/PCA128 sensitivity, direct pooled-source conformal, and attainable operating levels. The strict nested protocol uses $k\in\{1,2,4,8\}$, seeds from 0 to 4, $\alpha\in\{0.05,0.10,0.20\}$, and clean, Gaussian-noise, blur, brightness/contrast, and JPEG conditions. Each cell contains at most 120 evaluation images, with $\rho=0.01$, PCA64, $\delta=0.05$, and up to 20 candidate thresholds. Raw scores are first normalized using support-only statistics; condition-specific source views are then selected or aggregated as defined in Section~\ref{sec:cress-source-construction}. For every target category and seed, the eligible source categories are deterministically split using the frozen $0.50/0.25/0.25$ reference/proposal/certification allocation. This gives 7/3/4 categories for a 14-category within-MVTec source pool, 5/3/3 for an 11-category within-VisA pool, and 7/4/4 for the 15-category MVTec transfer pool. The same category partition is reused across conditions and $k$. The protocol evaluates MVTec and VisA within-dataset, MVTec-to-VisA, and MVTec-to-MPDD under matched-condition, clean-source, condition-agnostic, and mismatched negative-control source selection. The operational gate requires a nonzero category-certified threshold in at least 80\% of target cells, a target false-alarm rate (FAR) of at most $\alpha+0.02$, nonzero power, at least 80\% no-harm, and a positive lower bound on simultaneous power gain. Zero thresholds and failed cells are retained.

\subsection{Baselines and Ablations}
We compare the inherited PCA residual scorer with a controlled nearest-neighbor (NN) memory bank using the same frozen DINOv2 features and few-shot split. PCA64 and PCA128 retain 64 and 128 principal directions, respectively; PCA64 is the default reliability substrate. For the source-assisted study, direct pooled-source conformal is an uncertified baseline: it uses the complete selected source pool with $\tau=\alpha$, but performs no reference, proposal, or certification split and applies no UCB gate. The controlled memory-bank implementation is not an official reproduction of PatchCore or AnomalyDINO; it isolates scoring and storage trade-offs under a common protocol. Official AnomalyDINO code is also run on MVTec and VisA under its released protocol and reported as a separate accuracy reference. WinCLIP~\cite{winclip2023} is included only through published values because no official implementation is available. Official SubspaceAD representative runs serve as a novelty guardrail: the frozen DINOv2 subspace residual is inherited, and the contribution begins with the reliability analysis downstream of that score.

\subsection{Metrics}
Ranking metrics are distinguished from reliability metrics. AUROC and AP are calculated from the raw PCA residual score $s(x)$. Reliability is evaluated operationally by empirical FAR at fixed $\alpha$, detection power, alarm precision, the empirical cumulative distribution function of normal p-values, and the fraction of cells receiving a positive certified threshold. Storage is quantified as retained ranker state per category. Corruption views are controlled stress tests and are not treated as independent images or as a proxy for all deployment shifts.

\subsection{Reproducibility and Claim Rules}
We conduct all experiments on a single NVIDIA RTX 5090 graphics processing unit (GPU) using PyTorch 2.12.1. DINOv2 patch features are cached in 32-bit floating-point (fp32) format. For reproducibility, the frozen commit, DINOv2 source and weight hashes, environment specifications, dataset counts, per-image predictions, base-image identities, support manifests, corruption parameters, source partitions, candidate-threshold bounds, and Secure Hash Algorithm 256-bit (SHA-256) checksums are documented. Automated integrity checks encompass 1,500 MVTec, 1,200 VisA, and 600 MPDD view cells, identifying duplicate cells, missing support statistics, or overlap between support and test sets. Reliability methods cannot claim ranking improvement unless the raw score is altered; image-level and category-level certificates are not interchangeable.
